\section{AegisKV Design}
\label{sec:design}

AegisKV allocates finite KV storage to reduce workflow completion time
(WCT) while preserving normal cache reuse. Two observations connect
workflow-level priorities to the state needed for recovery.

\subsection{Observations}
\label{sec:design-insights}

\paragraph{1. Workflow dependencies determine recovery urgency.}
Figure~\ref{fig:motivation-recovery-exposure} shows similar reconstruction
costs with different effects on WCT: slack absorbs some branches' delays,
while a critical branch exposes its full pause. Thus request recovery
time alone is insufficient to rank protection priorities; the request's
position in the remaining workflow also matters.

\paragraph{2. Historical KV accesses concentrate on a subset of blocks.}
Figure~\ref{fig:sparse-releases} surveys long-context releases with block,
token or compressed-entry selection. Figure~\ref{fig:design-insights}
then examines native access traces: approximately 20\% of historical
blocks account for 87.64\%, 74.76\% and 56.54\% of accesses across the
three models. These results distinguish the complete historical state
from the subset frequently used by computation. The survey does not
establish a monotonic increase in sparsity, and retrospective concentration
does not establish future-query prediction accuracy.

\begin{figure}[t]
 \centering
 \includegraphics[width=\columnwidth]{figures/sparse-models.pdf}
 \caption{Selected sparse/hybrid model releases and reported context
 lengths. IDs: 1 MiniCPM4~\cite{minicpm4model};
 2 MiniCPM4.1~\cite{openbmb2025minicpm41};
 3 DeepSeek-V3.2-Exp~\cite{deepseek2025v32exp};
 4 DeepSeek-V3.2~\cite{deepseek2025v32model};
 5 GLM-5~\cite{zai2026glm5model};
 6 MiniCPM-SALA~\cite{openbmb2026sala};
 7 DeepSeek-V4 preview~\cite{deepseek2026v4model};
 8 MiniMax-M3~\cite{minimax2026m3};
 9 GLM-5.2~\cite{zai2026glm52};
 10 LongCat-Flash-Lite-Sparse~\cite{longcat2026sparse}.
 Dates are announcements, or first official repository dates for
 IDs 5, 9 and 10. Markers identify the sparse branch's selection unit,
 including SALA's block branch and V4's compressed sparse branch;
 they do not imply uniform attention across layers.}
 \label{fig:sparse-releases}
 \Description{Ten selected sparse or hybrid attention releases between
 June 2025 and July 2026, with reported contexts from 32K to 1M.
 Circles, squares and diamonds distinguish block, token and compressed
 entry selection. This is a literature survey, not a sparsity-rate curve.}
\end{figure}

\begin{figure}[t]
 \centering
 \includegraphics[width=\columnwidth]{figures/kv-access-concentration.pdf}
 \caption{Native KV access concentration for NOSA-8B~\cite{huang2026nosav2},
 NSA-8B (MiniCPM-4-8B-NSA)~\cite{fasa2026minicpm4nsa}, and
 MiniCPM4.1-8B~\cite{openbmb2025minicpm41}. We use the final 512 queries
 of matched 32,704-token prefixes from three PG19 texts. Each query
 counts a historical 64-token block once across heads. Rank the 503
 blocks within each layer/window; average access shares over all
 32 layers, then texts. Bands show text min--max; the diagonal denotes
 uniform access. Compressed/selector state is excluded. This is
 retrospective concentration, not prediction accuracy or recovery speedup.}
 \label{fig:design-insights}
 \Description{Cumulative historical block access curves for NOSA-8B,
 NSA-8B and MiniCPM4.1-8B, averaged over every layer and three texts.
 About 20 percent of blocks account for 87.64, 74.76 and 56.54 percent
 of accesses, respectively; all three curves exceed the uniform baseline.}
\end{figure}

Together, these observations guide which states to protect and which
parts to keep immediately accessible.

\subsection{System Overview}
\label{sec:system-overview}

AegisKV combines workflow priorities with a three-tier KV hierarchy
(Figure~\ref{fig:system-overview}). Peer HBM keeps selected
layers and hot blocks on another data-parallel (DP) replica for fast
recovery; shared DRAM and storage retain the remaining state for reuse
and recovery. Together, the three tiers cover each committed checkpoint,
with one managed owner per object after transfers finish. This division lets
the system retain long histories while spending scarce HBM on the state
whose faster access is most useful to workflow completion.

The Application supplies requests and workflow progress to the Controller.
\textbf{(1)} Workflow Tracker estimates how much delay each Agent can
absorb, while \textbf{(2)} Recovery-aware Router selects its execution and
takeover DPs. Using this urgency and location information, \textbf{(3)} KV
Master decides where to retain each KV object and which hot blocks merit
Peer space. Within the Engine, \textbf{(4)} Adapter adjusts the available
Peer capacity as serving demand changes. Connector executes Master's
transfer plans and reports their outcomes. If a worker fails, Router redirects its requests to the
takeover DP, where Connector combines Peer hits with lower-tier reads to
resume execution.

\begin{figure}[t]
 \centering
 \includegraphics[width=\columnwidth]{figures/system-overview.pdf}
 \caption{AegisKV architecture. Callouts (1)--(4) match the text;
 DP1 mirrors DP0's Adapter/Connector; the Master--Connector exchange is
 shown on DP1: plans flow down and transfer status returns up.
 DP0 executes; DP1 hosts selected Peer objects and takes over. Blue solid links
 carry KV, dark dashed links carry control, and red dashes denote retry;
 orange blocks are unsealed tails. The remaining checkpoint objects
 reside in shared DRAM/storage outside both DP failure boundaries.}
 \label{fig:system-overview}
 \Description{Compact Application, Controller and Engine layers with
 numbered components, symmetric DP Local/Peer HBM and shared lower tiers.}
\end{figure}

\subsection{Workflow Tracker: Remaining Time and Slack}
\label{sec:workflow-protection}

To make workflow priorities available online, Tracker estimates remaining
execution before deriving slack. It maintains a residual DAG whose nodes
are Agent segments between exposed fork/join boundaries or completion.
A segment can span several LLM calls and tool waits. Child execution
appears on child paths, not again in the parent; unknown future forks
are excluded.

A GBDT predicts segment seconds remaining, $d$, from causal features: role/model,
current input length, elapsed time, completed-call count and previous
calls' lengths, timings and context growth. Missing history is masked;
training, calibration and evaluation split whole workflows. Labels end
at the segment boundary, and future call outcomes are never features.

Algorithm~\ref{alg:workflow-tracker} takes residual DAG $G$ and features
$X$. Forward/backward critical-path passes~\cite{kelley1959cpm} compute
earliest finishes $EF$ and latest finishes $LF$ relative to the predicted
workflow finish. Their difference is nominal slack. The guard $g$ is an
upper empirical quantile of positive slack overestimation on calibration
workflows; subtracting it gives advisory slack $s$, not a per-request
guarantee. Requests in one segment share its margin.

\par\medskip\noindent\begin{minipage}{\columnwidth}
\raggedright
 \refstepcounter{designalgorithm}\label{alg:workflow-tracker}
 \textbf{Algorithm \thedesignalgorithm: Remaining time and slack}
 \par\smallskip
 \begin{tabularx}{\columnwidth}{@{}r@{\quad}X@{}}
 \toprule
 1 & $(G,X)\gets\textsc{SnapshotProgress}()$ \\
 2 & $\nu\gets\textsc{NextVersion}()$ \\
 3 & \textbf{if not} \textsc{Supported}$(G,X)$ \textbf{then} \\
 4 & \quad \textsc{PublishInvalid}$(\nu)$; \textbf{return} \\
 5 & $d\gets\max(0,\textsc{GBDT}(X));\quad g\gets\textsc{Guard}()$ \\
 6 & $EF\gets\textsc{Forward}(G,d)$ \\
 7 & $LF\gets\textsc{Backward}(G,d,\max_v EF_v)$ \\
 8 & \textbf{for} $v\in V(G)$ \textbf{do} \\
 9 & \quad $s_v\gets\max(0,LF_v-EF_v-g_v)$ \\
 10 & \textsc{Publish}$(\nu,G,d,g,s)$ \\
 \bottomrule
 \end{tabularx}
\end{minipage}\par\medskip

\textsc{Supported} checks inputs, dependencies, calibration and freshness;
invalid estimates grant no positive slack credit. Progress events and
periodic refresh publish versioned durations $d$, guards and slack to
Master, invalidating older placement costs. The DAG passes cost
$O(|V|+|E|)$; prediction quality is evaluated in
Section~\ref{sec:eval-prediction}.

\subsection{Recovery-aware Router}
\label{sec:recovery-router}

Router establishes where normal and resumed execution run. Normal
dispatch balances load and compatible prefix locality.
Figure~\ref{fig:agent-prefix-reuse} reports 97.90\% potential same-agent
token coverage and 0.37\% additional cross-agent coverage. Agent affinity
is a preference: shared prefixes can still help another Agent's first call.

\begin{figure}[t]
 \centering
 \includegraphics[width=\columnwidth]{figures/agent-prefix-reuse.pdf}
 \caption{Potential prefix reuse over 98,827 requests, 2,090 agents and
 393 CC/Weka workflows. Coverage uses earlier completed requests with
 the same recorded model and workflow. Same/other-agent coverage
 overlaps; only additional cross-agent coverage extends same-agent
 reuse. These are token-prefix opportunities, not runtime cache hits.}
 \label{fig:agent-prefix-reuse}
 \Description{Same-agent prefixes dominate aggregate reuse; cross-agent
 history can also help first subagent calls.}
\end{figure}

At arrival, Router records a versioned takeover binding to a DP with a
compatible layout outside the source failure domain. Its corresponding
tensor-parallel (TP) shards can own selected Peer objects. On confirmed
failure, Router fences the old epoch and activates that target subject
to workspace and admission constraints. A new target cannot count the
old target's Peer objects as local: rebinding retrieves reachable objects
or uses an older checkpoint/replay fallback if their sole owner was lost.
This binding gives Master concrete access paths and target load for
placement evaluation.

\subsection{KV Master: Workflow-aware Tiered Placement}
\label{sec:kv-master}

Master turns Tracker's workflow estimates and Router's takeover binding
into four decisions: which requests benefit from retention, which KV
objects carry that value, where to place them, and how to revise the
placement as demand changes.

\subsubsection{Workflow-aware Request Value}
\label{sec:placement-value}

A request's value comes from both normal reuse and failure recovery.
Retaining a prefix avoids reconstruction on a later call; faster recovery
also reduces WCT when its delay exceeds the request's guarded slack.
A large slack can absorb a recovery pause without making repeated
reconstruction free: those GPU cycles still compete with other requests.
Master therefore includes both benefits, using visible reuse history and
the current residual workflow. An existing local hit earns no additional
retention benefit.

Slack alone is insufficient to rank requests. The amount of work saved
and the capacity required also matter (Figure~\ref{fig:retention-value}).
Moreover, a source-worker failure can interrupt several branches at once;
accelerating only one of two tied branches may not advance their join.
Master ultimately compares complete placements $P$ and $P'$ through
\begin{equation}
 \Delta J=\sum_e p_e\sum_w[T_w(P',e)-T_w(P,e)].
 \label{eq:value}
\end{equation}
Here $T_w$ is predicted remaining WCT. The scenario table covers normal
continuation and source-domain faults at specified offsets, with
probabilities summing to one. Both plans use the same queues, finite
workspace and service profiles. Transfer interference enters request
execution times; the residual DAG propagates them through joins using
the maximum. Tracker's guard is applied once per interrupted segment.
Invalid workflow estimates fall back to request delays without positive
slack credit; invalid resource profiles pause optional moves.

\begin{figure}[t]
 \centering
 \includegraphics[width=\columnwidth]{figures/retention-value.pdf}
 \caption{Slack alone does not rank retention value. In this isolated-fault
 example, A uses two capacity units to reduce recovery from 4 to 2\,s;
 B uses one to reduce 18 to 8\,s. A at slack 0 saves 1\,s/unit, whereas
 B at slack 6 saves 10\,s/unit. Normal-use value is zero. These constructed
 profiles motivate cost/size comparison, not a policy-speedup claim.}
 \label{fig:retention-value}
 \Description{A higher-slack object can save more workflow time per byte.}
\end{figure}

\subsubsection{Prefix and Block Value}
\label{sec:prefix-value}

Request value must reach the prefix that makes reuse possible, rather
than accumulate on the newest block. For each Agent, Master starts from
a valid fallback and extends the recoverable prefix in token order.
Each increment includes the new token extent across all layers, TP shards
and required selector state. Its grant combines normal-use savings and
recovery delay removed beyond guarded slack, using the same horizon and
scenarios as Equation~\eqref{eq:value}. The previous restore path remains
an option, so increments are nonnegative. Restore times of 12, 7 and 5\,s
at successive frontiers, with slack 6\,s, yield recovery-value increments
of 5 and 1\,s before probability weighting. The final block does not
receive the entire 6\,s benefit.

Each increment is distributed over its newly required physical objects
in proportion to bytes. Distinct Agents contribute separate shares to a
shared prefix; repeated references by one Agent do not create new grants.
The physical object is charged once. Thus earlier prefixes can accumulate
value from more users, while later state retains its incremental benefit.
This allocation is a capacity-screening rule, not a claim that each layer
has independent recovery value. Grants are recomputed relative to the
available restore path; removing a dependency changes the value of the
remaining plan. Their local sum does not replace the joint WCT check.

Retention and immediate accessibility also need different priorities.
Keeping a recoverable prefix requires its full dependencies. Deciding
which of those objects deserve Peer HBM instead depends on when recovery
will access them and whether a lower-tier transfer would cause a wait.

\subsubsection{Three-Tier Placement and Peer Acceleration}
\label{sec:sparse-placement}

AegisKV uses exclusive placement across its managed tiers. A sealed
layer/token-block/TP object has one owning tier: Peer HBM, shared DRAM or
shared storage. A checkpoint is complete when the union of these objects
covers its prefix, including native auxiliary state. This managed state
is separate from the execution GPU's working KV. Moving an object to Peer
releases its lower-tier copy after transfer completes; capacity is not
spent on a permanent duplicate (Figure~\ref{fig:sparse-recovery}).

For request $r$, $L_r$ is the Peer prefix depth and $K_{r,l}$ counts the
selected token blocks in layer $l\le L_r$. Mandatory native dependencies
come first; other blocks are ranked by selection frequency over completed
queries, counted once per physical block per query across heads. Ties
favor recency, then block ID. A budget below the mandatory set is
infeasible. Selection units, engine blocks and kernel tiles need not
coincide; 64 tokens is a measurement unit here, not a universal requirement.
Objects outside the selected sets remain in DRAM or storage, and demand
misses retrieve their actual owners without changing native selections.

\begin{figure}[t]
 \centering
 \includegraphics[width=\columnwidth]{figures/sparse-recovery.pdf}
 \caption{One checkpoint distributed across three exclusive tiers.
 Each cell denotes a layer/block group, including its TP objects.
 Peer covers $L=2$ with $K_1=2$ and $K_2=1$; DRAM and storage hold the
 remaining cells. Layer 1 demands B0/B1/B2: B0 is local, B1 needs H2D,
 and B2 needs a storage read followed by H2D. The dashed tail is not yet
 committed. Demotion transfers ownership down the hierarchy; it cannot
 discard Peer bytes on the assumption of an existing lower copy.}
 \label{fig:sparse-recovery}
 \Description{A three-layer four-block ownership map with Peer, DRAM
 and storage owners, plus the three corresponding recovery access paths.}
\end{figure}

The last complete checkpoint covers $M$ tokens; the acknowledged tokens
beyond it form replay tail $N$. Saved history avoids reconstruction, while
computation on the tail can hide transfers for later layers. Peer earns
space only for the additional WCT reduction from avoiding exposed
lower-tier reads/H2D, relative to placing the same objects below it.
This benefit depends on $M$, $N$, hot-set coverage and recovery queues.
Peer uses direct revaluation, with no rental clock.

Figure~\ref{fig:design-pipeline} illustrates this distinction. Moving hot
sets from DRAM to Peer removes about 24\,ms of exposed wait while retaining
the same 1,024\,MiB of unique managed KV. The useful seven-layer depth is
specific to that history, tail and profile.

\begin{figure}[t]
 \centering
 \includegraphics[width=\columnwidth]{figures/recovery-pipeline.pdf}
 \caption{Analytical layered recovery at equal unique KV coverage.
 MiniCPM4-8B has 32 layers, 32K history and a 64-token tail. Replay is
 calibrated at 2.870\,ms/layer; assumed H2D bandwidth is 6.25\,GB/s
 (5.369\,ms/full layer). Trace-derived mean preload and preload-plus-miss
 fractions are 0.6584/0.7748. (a)--(c) hold 1,024\,MiB in DRAM; (d) moves
 147.5\,MiB to Peer in the first seven layers, leaving 876.5\,MiB in DRAM.
 Storage is inactive in this comparison. Numbers are tail completion
 and exposed wait; cold KV follows replay. Placement is already complete,
 and the ideal preemptible stream omits movement, finite-workspace and
 foreground-queue costs. These are analytical schedules, not end-to-end
 recovery measurements.}
 \label{fig:design-pipeline}
 \Description{Four recovery schedules with identical replay computation
 and unique KV coverage, reducing exposed wait from 172 to 83, 44 and
 20 milliseconds through overlap, sparse preloading and Peer placement.}
\end{figure}

Master enumerates bounded block-count/fraction templates and depths from
zero to model depth, retaining the current placement as an option. Each
fault scenario uses current $M$ and its predicted $N$, including commit
lag; unsupported sparse profiles use full demand. Only computation that
must still execute hides transfer. Neither $L$ nor $K$ must grow with
every token: block identities can change at fixed size, and $L=0$ is valid.

Lower-tier placement and Peer selection are evaluated together. New cold
state may enter DRAM or storage directly. Peer demotion moves objects to
DRAM; DRAM demotion moves them to storage. A full destination may require
a cascade before Peer memory becomes reusable. Final storage eviction
must publish an older valid checkpoint or input/log replay fallback for
every affected manifest. A shared object has one owner even when Agents
have different takeover targets: only its host DP gets a local Peer hit;
other access paths and the cost of moving it enter the plan evaluation.
The protected source failure must leave all advertised owners reachable;
loss of a sole owner requires fallback, not an assumed hidden replica.

\subsubsection{Renewable Credit and Online Management}
\label{sec:placement-plan}

Repeated calls and shared prefixes create demand that outlives one request.
A one-time score misses new users, while renewing credit at every
checkpoint can preserve stale state indefinitely. Master uses renewable
credit in DRAM and storage, inspired by Landlord and
GreedyDual-Size~\cite{young2002landlord,cao1997greedydual}, while revaluing
Peer acceleration directly as described above.

For lower pool $j$, object $x$ has size $b_x$, Agent grants $g_{a,x,j}$ and
paid amounts $d_{a,x,j}$. With cumulative pressure rent $\rho_j$, its
remaining credit and eviction key are
\begin{equation}
 C_{x,j}=\sum_a[g_{a,x,j}-d_{a,x,j}]^+,
 \qquad H_{x,j}=\rho_j+C_{x,j}/b_x.
 \label{eq:rental-credit}
\end{equation}
Each pool starts at zero. A rent increase $\delta$ charges every committed
resident object $\min(C_{x,j},b_x\delta)$, proportionally across its positive
Agent shares. Pinned objects also age; removability is a separate property.
Rent measures capacity competition rather than elapsed time.

Actual reuse renews only the using Agent's share; a use ID prevents repeated
renewal within that use. Value changes refresh grants while preserving
paid amounts. Checkpoint growth and tier migration are not reuse.
Per-pool paid amounts survive tier round trips, but accrue only while
resident in that pool; a move through Peer cannot erase earlier rent.
A new use resets that Agent's paid amounts. Tool waits preserve forecast
interest; expired interest removes grants. Figure~\ref{fig:renewable-credit}
illustrates why renewal differs from revaluation.

\begin{figure}[t]
 \centering
 \includegraphics[width=\columnwidth]{figures/renewable-credit.pdf}
 \caption{Lower-tier credit for one shared, unit-size DRAM object.
 Initial Agent grants are 3 and 2. Rent of 2 leaves shares 1.8 and 1.2.
 Updating A's grant to 4 preserves its paid amount 1.2, leaving 2.8;
 an actual new use by A then renews its share to 4. B stays at 1.2.
 Credit screens capacity competition; it is neither measured WCT nor
 duplicated capacity. Peer residency uses no rent.}
 \label{fig:renewable-credit}
 \Description{A lower-tier ledger separating pressure aging, value-only
 refresh and actual-use renewal for two Agents sharing one object.}
\end{figure}

\paragraph{Admission and victim search.}
For DRAM/storage, Master refreshes $q\ge2$ low-key objects and $q$ rotating
alternatives within a common work budget. Victim bundles must release
actual bytes without breaking manifests. Releasing lower-tier objects
$Y$ raises that pool's trial clock to at least $\max_{x\in Y}H_{x,j}$.
Incoming bytes must retain positive credit after the same trial charge;
existing objects receive no duplicate grant. Peer victim candidates use
current recovery loss per byte, subject to the same legal-plan check.
Bounded bundles can also advance tied workflow branches together.

Ordinary competition respects running requests' logical Peer pins.
Request end/tool handoff unpins; self hot-set changes and required quota
relief can override logical pins, never physical read/copy leases.
Every plan reserves destination and temporary space, recursively including
lower-tier victims. Optional plans require positive incoming lower-tier
credit and $\Delta J<-\epsilon$. Mandatory capacity relief bypasses the
positive-credit gate and chooses the least costly examined legal plan.
If no safe cascade fits, admission waits; shrinking a quota does not
instantly free occupied HBM.

\paragraph{Closing the loop.}
Prefill offers newly sealed state; decode triggers review every configured,
block-aligned $Z$ tokens. Request completion/tool handoff offers a final
snapshot, including short rounds. Actual use, value changes, capacity
updates, failure and transfer completion also trigger Algorithm~\ref{alg:placement}.
Candidates cover prefix extension, promotion/demotion, valid final-copy
removal and Peer hot-set replacement. A budgeted refresh revisits stale
keys. If an extension cannot commit, $M$ stays unchanged; copy lag may
make $N>Z$.

\par\medskip\noindent\begin{minipage}{\columnwidth}
\raggedright
 \refstepcounter{designalgorithm}\label{alg:placement}
 \textbf{Algorithm \thedesignalgorithm: Online three-tier placement}
 \par\smallskip
 \begin{tabularx}{\columnwidth}{@{}r@{\quad}X@{}}
 \toprule
 1 & \textsc{Reconcile}$(e)$; \textsc{SettleLowerResidents}$(e)$ \\
 2 & \textbf{if} \textsc{NewUse}$(e)$ \textbf{then} \textsc{RenewAgent}$(e)$ \\
 3 & $S\gets\textsc{RevalueAndSnapshot}(e)$ \\
 4 & $h\gets\textsc{ShrinkRequired}(S)$; $(P^*,d^*)\gets(\varnothing,\infty)$ \\
 5 & \textbf{for} $c\in\textsc{Candidates}(S,h)$ \textbf{while} \textsc{WorkLeft}() \\
 6 & \quad \textbf{for} $P\in\textsc{LegalPlans}(S,c,h)$ \\
 7 & \qquad \textbf{if} $\neg h\land\neg\textsc{LowerRentOK}(P)$ \textbf{then continue} \\
 8 & \qquad $d\gets\Delta J(P;S)$ \\
 9 & \qquad \textbf{if} $d<d^*$ \textbf{then} $(P^*,d^*)\gets(P,d)$ \\
 10 & \textbf{if} $P^*=\varnothing$ \textbf{then} \textsc{Defer}$(e)$; \textbf{return} \\
 11 & \textbf{if} $\neg h\land d^*\ge-\epsilon$ \textbf{then return} \\
 12 & \textbf{if not} \textsc{ReserveIfCurrent}$(S,P^*)$ \textbf{then} \\
 13 & \quad \textsc{Requeue}$(e)$; \textbf{return} \\
 14 & \textsc{Submit}$(P^*)$ \quad\textit{completion/failure emits a new event} \\
 \bottomrule
 \end{tabularx}
\end{minipage}\par\medskip

Connector copies each object, confirms the destination, redirects its
location, then releases the source after readers drain. Both physical
copies consume capacity during this handoff. Versions, leases and pool
clocks are rechecked before reservation; rejected previews change neither
rent nor ownership. Trial rent becomes committed at the corresponding
capacity handoff, and admitted objects keep their residual trial credit.
An aborted copy before handoff preserves the old location; later failure
cannot undo completed moves or rent. Only complete dependency coverage advances the
checkpoint. This bounded search offers no optimality or inherited
competitive guarantee.

\subsection{Adapter: Guarded AIMD for Peer Capacity}
\label{sec:adapter}

Master selects content within a quota; Adapter adjusts that quota as
foreground inference demand changes. It uses pressure-driven multiplicative
decrease and stable-window additive increase~\cite{chiu1989aimd}.
The total Peer authorization has a physical ceiling
\begin{equation}
 C=\max\{0,\min(B_{\mathrm{cap}},
 B_{\mathrm{pool}}-B_{\mathrm{serve}}-B_{\mathrm{reserve}})\}.
 \label{eq:peer-ceiling}
\end{equation}
The pool excludes weights/fixed allocations; serving high watermark
excludes Peer, and reserve covers growth/workspace. Existing Peer bytes
are not subtracted twice. $B_{\rm cap}$ can be a configured pool fraction.

The current quota $B$ and stable-window count $z$ start at zero.
Pressure (allocation stalls, reserve violations or
memory-related admission waits) multiplies $B$ by $0<\gamma<1$. After $m$
stable windows below a lower wait threshold, deferred positive-value Peer
demand, assessed before quota filtering, permits an additive byte step
$\Delta$. Intermediate windows reset
$z$; missing telemetry disables growth. The ceiling clips every update;
emergencies may decrease within a window, but growth requires completed
stable windows. Parameters and window length are configurable.

\par\medskip\noindent\begin{minipage}{\columnwidth}
\raggedright
 \refstepcounter{designalgorithm}\label{alg:adapter}
 \textbf{Algorithm \thedesignalgorithm: Windowed Peer quota}
 \par\smallskip
 \begin{tabularx}{\columnwidth}{@{}r@{\quad}X@{}}
 \toprule
 1 & $C\gets\textsc{SafeCeiling}();\quad B'\gets B$ \\
 2 & \textbf{if} \textsc{Pressure}() \textbf{then} $(B',z)\gets(\gamma B,0)$ \\
 3 & \textbf{else if} \textsc{Stable}() \textbf{then} \\
 4 & \quad $z\gets z+1$ \\
 5 & \quad \textbf{if} $z\ge m$ \textbf{then} \\
 6 & \qquad \textbf{if} \textsc{PeerDemand}() \textbf{then} $B'\gets B+\Delta$ \\
 7 & \qquad $z\gets0$ \\
 8 & \textbf{else} $z\gets0$ \\
 9 & $B\gets\textsc{RoundDownPage}(\min(C,B'))$ \\
 10 & \textsc{PublishBudget}$(B)$; \textsc{EmitCapacityEvent}() \\
 \bottomrule
 \end{tabularx}
\end{minipage}\par\medskip

Adapter sets the total budget; Master chooses $L/K$ at fixed $Z$ and
other placement events. If occupancy plus reservations exceeds a reduced
quota, Master pauses Peer admission and selects objects to demote. Their
DRAM destinations may first require storage demotion. Serving can reuse
the memory only after these transfers and ongoing readers finish; the
new quota is an authorization target, not an instantaneous free-space claim.
